
\subsubsection{2.1 Token Log-Probability Baselines}

Token-level log-probabilities require no additional inference beyond the standard forward pass. Fadeeva et al. [2024] introduce the Claim Conditioned Probability (CCP), which uses mean log-probability while conditioning on the factual claim portion of a response, achieving AUROC 0.72--0.80 on TriviaQA and Natural Questions across seven LLMs and four languages. Kadavath et al. [2022] demonstrate that token log-probabilities correlate with model self-knowledge via the P(True) and P(IK) constructs, establishing the empirical basis for log-probability uncertainty signals. The present work extends this line by showing that mean aggregation --- the implicit choice in CCP --- is suboptimal for factual recall benchmarks by 5--12 AUROC points when minimum is used instead.

\subsubsection{2.2 Multi-Sample Consistency Methods}

Manakul et al. [2023] introduce SelfCheckGPT, which detects hallucination by measuring semantic consistency across N sampled outputs (N=20 for competitive performance), achieving AUROC 0.72--0.78 on sentence-level hallucination detection. SelfCheckGPT includes ad hoc token log-probability baselines without systematic ablation of aggregation function or directional predictions. Farquhar et al. [2024] introduce Semantic Entropy (SE), which clusters N sampled outputs by semantic equivalence via NLI and computes predictive entropy over clusters, achieving AUROC $\approx 0.79$ on TriviaQA with N=10 samples. SE's within-cluster sum is an implicit aggregation choice that is not ablated. The best single-pass result in this work (raw sum AUROC 0.895--0.896 on TriviaQA) exceeds published SE AUROC at one-tenth the inference cost, though cross-pipeline comparison warrants caution (see Section 6.3).

Recent extensions reduce SE's sampling overhead: Ciosek et al. [2025] achieve matching AUROC with 53\% fewer samples via Bayesian SE, and Nguyen et al. [2025] improve via pairwise semantic similarity (SNNE). These methods remain in the multi-sample regime. The present work is orthogonal: it addresses the single-pass regime and provides the first principled aggregation selection criterion within it.

\subsubsection{2.3 Supervised and Probe-Based Methods}

Shelmanov et al. [2025] train uncertainty quantification heads on attention maps to detect hallucination at claim level across Mistral, LLaMA, and Gemma, achieving state-of-the-art performance at the cost of supervised training per model. Moslonka et al. [2025] propose EPR, which extracts features from top-k logits and applies a learned detector. The present work is entirely unsupervised: no fine-tuning, no learned components, no labeled calibration data. \textit{Note: Moslonka et al. [2025] and other 2025 arXiv preprints cited in this work should be verified against published versions prior to final submission.}

\subsubsection{2.4 Hallucination Benchmarks and Taxonomy}

Lin et al. [2022] introduce TruthfulQA, a benchmark of 817 questions designed to elicit imitative-falsehood hallucinations. Farquhar et al. [2024] provide reproducible binary-label splits of TriviaQA from the \texttt{jlko/semantic\textbackslash \_uncertainty} repository. The taxonomy introduced here --- hallucination type determines token probability distribution shape, which determines optimal aggregation function --- is the first mechanistic classification that predicts optimal aggregation from hallucination type alone.

\subsubsection{2.5 Aggregation Function Comparison: The Open Gap}

Liu et al. [2025] survey uncertainty quantification and calibration for LLMs, identifying systematic comparison of token-level log-probability aggregation strategies as an open challenge. The present paper fills this gap through controlled ablation with mechanism-grounded directional predictions confirmed with bootstrap confidence intervals across two model families.

